% !TeX root = main.tex
\section{Introduction}
\label{sec:introduction}
\begingroup
\setlength{\emergencystretch}{1em}

Video recommendation on Web platforms depends on representations that connect visual content with user preferences.
Multimodal large language models (MLLMs) provide semantic information beyond item metadata: DWYS~\cite{denadai2025dwys} uses generated descriptions as recommendation inputs, while IISAN~\cite{fu2024iisan} adapts multimodal representations through side networks.
Resource constraints can favor compact MLLMs for content extraction.
This raises a representation learning question: how can a recommender make better use of the information that a compact MLLM already extracts?

We investigate explicit structure as a way to improve the recommendation utility of compact MLLM outputs.
A scene may contain entities and actions, yet their separate aggregation leaves participation roles implicit; likewise, aggregating scene vectors without action-level organization leaves continuity implicit.
The opportunity is to expose these relations to the recommender through structured representations of the available observations.

For example, a person may cut vegetables with a knife across several consecutive scenes.
The person, vegetables, and knife play different roles in the same action, while the repeated action links the scenes into a coherent segment.
This suggests contiguous action segments as aggregation units for video representation.
We examine this choice within each video; user interaction histories remain modeled by a sequential recommender.

We introduce HGVRec, a hierarchical framework centered on Action-Segment Hypergraph Aggregation (ASHA).
Building on the general aggregation framework of AllSet~\cite{chien2022allset}, ASHA treats scenes as nodes and maximal contiguous runs sharing a normalized action key as hyperedges.
Scenes with multiple actions can belong to overlapping segments.
Messages pass from scenes to segments and back, with segment extent incorporated into the representations and video readout.
Role-aware Entity--Action Fusion (REAF) supplies ASHA with scene features by fusing entity features into actions according to extracted actor, target, tool, and place assignments.
These explicit memberships and role links make the modeled structure inspectable.

We frame evaluation around structural modeling with a fixed compact MLLM, comparisons with scene sequence and attention encoders, and the use of more capable extractors.
As device chips improve or server AI becomes available, stronger MLLMs may become practical; HGVRec's modular extraction stage allows their recommendation benefits to be investigated within the same framework.
Our contributions and the empirical questions they motivate are:
\begin{itemize}
  \item \textbf{Structured representations from compact MLLMs.}
  We propose a framework that explicitly organizes extracted semantics into role-aware scenes and action segments to improve their utility for recommendation.
  % TODO (A): Verify the gain with the same compact MLLM, visual evidence,
  % recommendation backbone, and evaluation protocol before stating a result.
  \item \textbf{Action segments as aggregation units.}
  We formulate contiguous action grouping as a video representation unit and investigate its utility relative to chronological scene encoding.
  % TODO (B): Add a matched comparison using identical scene features with
  % chronological sequence encoding; report results only after verification.
  \item \textbf{Inspectable relational modeling.}
  ASHA and REAF define explicit scene memberships and role links, supporting structural inspection and comparison with attention-based scene aggregation.
  % TODO (C): Add a matched attention comparison and an actual segment/role
  % example. Structural inspection alone does not establish explanation fidelity.
  \item \textbf{Extension across MLLM capabilities.}
  The framework separates semantic extraction from recommendation learning, enabling evaluation of compact and more capable MLLMs under a common recommendation setup.
  % TODO (D): Add verified comparisons across extractors under the same
  % evaluation protocol; do not infer device latency or memory scalability.
\end{itemize}
\par
\endgroup

